\section{Discussion}
\label{sec:discussion}

\subsection{Main Findings}
\label{subsec:main_findings}

The experiments provide evidence that both temporal and spatial information
are relevant for traffic forecasting on METR-LA. The Temporal model already
provides a substantial improvement over persistence, showing that recent
sensor histories contain strong predictive information. Incorporating graph
message passing further improves forecasting accuracy, with the benefit
becoming more evident at longer forecasting horizons. Moreover, the
sensor-level analysis shows that this improvement is widespread across the
network, although a small number of sensors do not benefit from spatial
information.

The graph ablations further show that the way spatial information is
represented matters. Increasing the diffusion depth from $K=2$ to $K=3$
provides only limited additional benefits and introduces a less favorable
balance between error metrics, calibration, and computational cost. In
contrast, replacing the weighted adjacency matrix with a binary one produces
a clearer degradation in forecasting accuracy, particularly at longer
horizons. This suggests that the original edge weights provide useful
information beyond graph connectivity alone.

The probabilistic analysis reveals a different aspect of model performance.
The Graph-Temporal model produces sharper intervals and lower quantile losses
than the Temporal model, but its empirical coverage is slightly below the
nominal 80\% level. More importantly, this reliability degradation becomes
stronger under distribution shift. During high-congestion periods, both point
forecasting errors and prediction interval widths increase, while coverage
decreases at longer horizons. The model therefore recognizes part of the
increased uncertainty associated with congestion, but does not fully capture
the shifted predictive distribution.

Finally, the mitigation experiments show that post-hoc calibration can
substantially improve uncertainty reliability. CQR brings overall empirical
coverage close to the nominal level without modifying the median forecasts or
retraining the model. However, residual undercoverage remains under
high-congestion conditions, particularly at the longest horizon. This result
indicates that good marginal calibration does not necessarily imply equally
reliable uncertainty estimates under specific shifted conditions and
motivates mitigation strategies that explicitly account for rare traffic
regimes during training.


\subsection{Limitations}
\label{subsec:limitations}

Several limitations should be considered when interpreting the results.
First, the experiments are conducted exclusively on METR-LA and therefore
reflect the spatial structure, traffic dynamics, and missing-data patterns of
a single road network. The conclusions cannot necessarily be generalized to
other cities or sensor networks without additional evaluation.

Second, the definition of distribution shift considered in this work focuses
specifically on high-congestion periods. Although this represents a relevant
and challenging traffic regime, other forms of shift may affect forecasting
models differently, including sensor failures, changes in network topology,
or temporal changes in traffic patterns.

Third, the Graph-Temporal architecture is intentionally compact and explores
only a limited set of spatial configurations. The diffusion-depth and
adjacency ablations provide evidence about the role of graph information, but
they do not constitute an exhaustive exploration of graph neural network
architectures or graph construction strategies.

Finally, CQR is calibrated globally for each forecasting horizon using the
validation distribution. While this substantially improves marginal
coverage, it does not explicitly account for traffic regime and therefore
cannot guarantee the same calibration quality within the high-congestion
subset. In addition, the temporal and spatial dependence present in traffic
data makes the standard exchangeability assumptions underlying conformal
prediction less directly applicable.


\subsection{Future Improvements}
\label{subsec:future_improvements}

Several extensions could address these limitations. A first direction would
be to evaluate the proposed congestion-aware training objective, assigning
greater importance to rare congested targets during optimization and
measuring whether this improves both point forecasts and uncertainty
estimation without substantially degrading performance under normal
conditions.

The robustness analysis could also be extended to additional distribution
shifts, such as partially masked sensor histories or simulated sensor
failures. This would provide a broader characterization of the dependence of
the Graph-Temporal model on spatial information and of its behavior when
parts of the network become unavailable.

Finally, future work could explore more expressive spatial and temporal
architectures, alternative graph construction strategies, and dynamic graphs
whose relationships evolve with traffic conditions. Evaluation on additional
traffic datasets would also be necessary to determine whether the conclusions
observed on METR-LA generalize to different road networks and traffic
regimes.